\section{A common denominator}\label{sec:arithmetic}

Put $d_n=\lcm(1,\ldots,n)$ and $\delta_n=4d_n$. For a prime $p$,
$v_p(x)$ is the exponent of $p$ in a nonzero rational number $x$, with
$v_p(0)=+\infty$. A rational number is \emph{integral at $p$} if its
valuation is nonnegative. We bound the denominators of both the beta
coefficients and the rational correction terms.

\begin{proposition}\label{prop:denominators}
For $1\leq i\leq364$, $0\leq h\leq n$ and $b=1,2,3$,
\begin{equation}\label{eq:denominators}
 \delta_n^{365-i}a_{i,h}(n)\in\Z,\qquad
 \delta_n^{365-i}\rho_i(n)\in\Z,\qquad
 \delta_n^{365}q_b(n)\in\Z.
\end{equation}
\end{proposition}

\subsection{Taylor coefficients of the numerator}

Define
\[
 F_\theta(t)=\frac{A_\theta^n}{(n!)^{18}}
          \prod_{j=0}^{18n-1}(t-7n+j+\theta),
 \qquad
 A_1=1,\quad A_{1/2}=2^{36},\quad A_{1/3}=A_{2/3}=3^{27}.
\]
For every integer $u$ and every $k\geq0$, we claim
\begin{equation}\label{eq:taylor-integral}
 d_n^k[t^k]F_\theta(t+u)\in\Z.
\end{equation}
Here $[t^k]$ denotes the coefficient of $t^k$. Write $\theta=a/b$ in
lowest terms. Multiplying each linear factor of $F_\theta(t+u)$ by $b$
gives integer constant terms in a progression with step $b$.
Fix a prime $q\nmid b$ and split this progression into $18$ blocks of
length $n$. For every $q^v\leq n$, at least
$18\lfloor n/q^v\rfloor$ entries are divisible by $q^v$.
To form a term in the coefficient of $t^k$, we choose the linear term
from $k$ factors and constant terms from the rest. Omitting those $k$ constants
loses at most $k$ contributions at each prime-power level. Before
division by $(n!)^{18}$, the resulting valuation is at least
\[
 18v_q(n!)-k\lfloor\log_q n\rfloor.
\]
Multiplication by $d_n^k$ compensates for this loss. Any term that
still contains a zero constant is zero and needs no estimate.

The factor $A_\theta^n$ compensates for primes dividing $b$.
For $b=2$, after clearing the $18n$ factor denominators, $2^{18n}$
remains; this covers $18v_2(n!)\leq18n$.
For $b=3$, $3^{9n}$ remains, covering $18v_3(n!)\leq9n$.
Coefficient extraction introduces only further integer factors.
Thus \eqref{eq:taylor-integral} holds at every prime.

Let
\[
 F(t)=(t-7n)\prod_{\theta\in\Theta}F_\theta(t)
     =\frac{K^nP_n(t)}{(n!)^{72}}.
\]
Expansion at $-4h$ gives
\begin{align}
 a_{i,h}(n)
  &=(n+1)^{365}(-1)^{365h}\binom nh^{365}\notag\\[-2pt]
  &\quad\times[t^{365-i}]
   \left\{F(t-4h)
      \prod_{\substack{0\leq k\leq n\\k\ne h}}
       \left(1+\frac{t}{4(k-h)}\right)^{-365}\right\}.
       \label{eq:local-coefficient}
\end{align}
Each $4(k-h)$ divides $\delta_n$. The binomial series for the remaining
factors, together with \eqref{eq:taylor-integral}, proves the first
assertion of \eqref{eq:denominators}; summing over $h$ proves the second.
The cancellation of the constant powers of $4$ uses the factor
$4^{365n}$ in $C_n$; the same factor also contributes to the decay
estimate in Section~\ref{sec:decay}.

\subsection{Clearing the rational correction terms}

The denominators $j$ in \eqref{eq:qb} can exceed $n$, so multiplication
by $\delta_n^{365}$ need not make each summand integral. We group the
terms with the same $h,j$ by summing over the pole orders $i$.
At any fixed prime, at most one pole group can fail to be integral after
scaling. At a zero of $R_n$, the sum of all pole groups is zero.
The possible exceptional group is therefore the negative of a sum of
integral groups, and is integral as well.

Fix an odd $j$ with $1\leq j\leq4bh$. The point $t_0=j/b-4h$ lies in
$(-4n,0)$ and is a numerator zero: its reduced denominator divides $b$,
and $\Theta$ supplies the required fraction class. It is not a pole,
since $t_0=-4k$ would imply $j=4b(h-k)$, impossible for odd $j$.
Evaluating $R_n(t_0)=0$ yields
\begin{equation}\label{eq:grouped-zero}
 \sum_i\frac{a_{i,h}(n)}{(j/b)^i}
 =-\sum_{k\ne h}\sum_i
            \frac{a_{i,k}(n)}{(j/b+4(k-h))^i}.
\end{equation}

Fix a prime $q$, put $v=v_q(\delta_n)$, and set $x_k=j/b+4(k-h)$.
At most one $x_k$ can have $v_q(x_k)>v$: two such numbers would have
a difference with valuation greater than $v$, whereas
\[
 x_k-x_l=4(k-l)\mid\delta_n\qquad(k\ne l).
\]
The coefficient bound already proved implies, for every nonzero rational $x$,
\begin{equation}\label{eq:valuation-group}
 v_q\left(\frac{\delta_n^{365}a_{i,h}(n)}{x^i}\right)
       \geq i\bigl(v-v_q(x)\bigr).
\end{equation}
If $v_q(x_h)\leq v$, this inequality makes every term on the left of
\eqref{eq:grouped-zero} integral at $q$ after multiplication by
$\delta_n^{365}$. If $v_q(x_h)>v$, every other $x_k$ satisfies
$v_q(x_k)\leq v$. The scaled right side is then
integral at $q$, and the identity proves the same for the left side.
Thus every scaled group is integral at every prime. Summing the groups
in \eqref{eq:qb} proves the last assertion of
Proposition~\ref{prop:denominators}.
